\begin{solapartat}
\punts{0,4} Primer determinarem la direcció i el sentit del camp magnètic creat pel fil (1). Sabem que les línies de camp són circumferències en el pla $xy$ centrades en el fil, i el sentit es determina amb la mà dreta posant el dit polze en el mateix sentit que el corrent, com s'indica a la figura.

\figura[0.2]{sol_fig1.png}

Com que el camp magnètic és tangent a les línies i té el mateix sentit, el camp magnètic creat pel fil (1) té la direcció i el sentit indicats a la figura. Noteu que, com que és tangent a la circumferència, és perpendicular al radi. Aquest camp magnètic té dues components segons els eixos $x$ i $y$, o el que és equivalent, forma un cert angle $\alpha$ amb l'eix $y$.

\figura[0.35]{sol_fig2.png}

\punts{0,4} Seguint el mateix raonament anterior i tornant a aplicar la regla de la mà dreta, es veu immediatament que el sentit i la direcció del camp magnètic creat pel fil 2 al punt $P$ són els indicats a la figura adjunta.

\figura[0.35]{sol_fig3.png}

\punts{0,2} Els mòduls dels camps $B_1$ i $B_2$ són iguals, atès que el corrent que circula pels dos fils és el mateix i el punt és equidistant als dos fils.

\punts{0,25} A més, com que el punt $P$ és equidistant als dos cables, llavors per simetria els vectors $\vec{B}_1$ i $\vec{B}_2$ formen el mateix angle $\alpha$ respecte a l'eix $y$, i com que $|B_y| = B\cos(\alpha)$, llavors $|B_{y,1}| = |B_{y,2}|$. A més, de l'anàlisi anterior de la direcció i el sentit dels camps $\vec{B}_1$ i $\vec{B}_2$, tenim que $B_{y,2} = -B_{y,1}$. En canvi, les dues components en l'eix de les $x$ són positives. Per tant, quan fem la suma vectorial, les components en la direcció $y$ es cancel·len mútuament, de manera que només ens queda la suma de les components en la direcció $x$:

\figura[0.4]{sol_fig4.png}
\end{solapartat}

\begin{solapartat}
\punts{0,4} Primer calcularem la magnitud del camp magnètic:
\[ B_1 = \frac{\mu_0 I}{2\pi r} = \frac{4\pi\times 10^{-7}\cdot 1,5}{2\pi\cdot 0,16} = 1,875\times 10^{-6}\ \mathrm{T} \]

\punts{0,25} Si apliquem la regla de la mà dreta, el camp magnètic va dirigit segons l'eix $x$ i en la direcció positiva:
\[ \vec{B}_1 = 1,875\times 10^{-6}\,\vec{\imath}\ \mathrm{T} \]

Seguidament, calcularem el mòdul de la força:
\[ \vec{F}_m = I(\vec{l}\times \vec{B}_1) \]

\punts{0,35} El camp magnètic va dirigit segons l'eix $x$, mentre que $\vec{l}$ va dirigit segons l'eix $z$; per tant, són perpendiculars i el mòdul de la força per unitat de longitud serà:
\[ F_m = I l B_1 \;\Rightarrow\; \frac{F_m}{l} = I B_1 = 2,81\times 10^{-6}\ \mathrm{N/m} \]

També es considerarà correcte que es doni com a resultat la força aplicada sobre un metre de fil:
\[ F_m = I l B_1 = 1,5\ \mathrm{A}\cdot 1\ \mathrm{m}\cdot 1,875\times 10^{-6}\ \mathrm{T} = 2,81\times 10^{-6}\ \mathrm{N} \]

\punts{0,25} Finalment, la direcció de $\vec{F}_m$ és perpendicular tant a $\vec{l}$, eix $z$, com a $\vec{B}_1$, eix $x$; per tant, $\vec{F}_m$ va dirigida segons l'eix $y$. El sentit el determinem a partir de la regla de la mà dreta:
\[ \frac{\vec{F}_m}{l} = 2,81\times 10^{-6}\,\vec{\jmath}\ \mathrm{N/m} \]

\figura[0.6]{sol_fig5.png}
\end{solapartat}
